% source: https://github.com/pfradin/luadraw/discussions/289#discussioncomment-17611518 (pfradin, block 1)
\documentclass[border=3pt]{standalone}% compile with lualatex only
\usepackage[svgnames]{xcolor}
\usepackage[3d]{luadraw}%https://github.com/pfradin/luadraw
\usepackage{fourier-otf}
\begin{document}
\begin{luadraw}{name=cylinder-cone-sphere}
local ld = luadraw
require 'luadraw_spherical'

local pt3d = ld.pt3d
local M, Mc, O, vecK = pt3d.M, pt3d.Mc, pt3d.Origin,  pt3d.vecK
local sin, cos, pi, sqrt = math.sin, math.cos, math.pi, math.sqrt
local Rsphere = 5
local Rcone = 2 * sqrt(2) * Rsphere / 3
local hcone = 4 / 3 * Rsphere
local B = M(0, 0, -Rsphere / 3) 
local T = M(0, 0, hcone - Rsphere / 3)
local Rcylinder = 2 * Rcone / 3
local hcylinder = hcone / 3
local H_cyl = M(0, 0, -Rsphere / 3 + hcylinder) 
local g = ld.graph3d:new{window={-Rsphere-1,Rsphere+1,-Rsphere-1,Rsphere+1},viewdir={70,60},size={12,12}}
ld.Hiddenlinestyle = "dashed" -- after the declaration of graph g

-- points of tangency for the cone
local I, J = pt3d.vecI, pt3d.vecJ -- *(B, I, J)* is an orthogonal coordinate system for the base of the cone.
local f = function(t) 
    local R = Rcone
    local H = T
-- M = B+R*cos(t)*I+R*sin(t)*J is a point on the circular base.
-- f returns the determinant of the vectors g.Normal, M-H, and d(M-H)/dt
    return pt3d.det( g.Normal, B+R*cos(t)*I+R*sin(t)*J-H, -R*sin(t)*I+R*cos(t)*J ) 
end
-- The points of tangency are the points on the circle where the function f vanishes.
local t = ld.solve(f, -pi, pi)
local C, D = B+Rcone*(cos(t[1])*I+sin(t[1])*J), B+Rcone*(cos(t[2])*I+sin(t[2])*J)
-- 

g:Dcylinder(B, Rcylinder, H_cyl, {edgestyle="dashed", edgecolor="red"}) 
--g:Dcone(B, Rcone, T, {edgestyle="dashed", edgecolor="green"})
g:Dpolyline3d( {C,T,D}, "green,dashed")
g:Define_sphere({radius=Rsphere, center=O, mode=ld.mBorder, color="", edgecolor="blue"})
g:DScircle({B, vecK}, {hidden=true,color="green"}) 
g:Dspherical()
g:Show()
\end{luadraw}
\end{document}
